\documentclass[10pt,a4paper]{amsart}
\usepackage[margin=19mm]{geometry}
\usepackage{amsmath,amssymb,mathtools}
\usepackage[scr=rsfs,cal=euler]{mathalfa}
\usepackage{xfrac}
\usepackage[unicode,hidelinks,bookmarks=false]{hyperref}
\newcommand{\ie}{i.e.\ }
\newcommand{\cf}{cf.\ }
\newcommand{\resp}{respectively\ }
\newcommand{\Subsection}{\subsection}
\providecommand{\nobreakdash}{\nobreak\hbox{-}}
\setlength{\emergencystretch}{2em}
\multlinegap0pt
\usepackage{bm}
\usepackage{amssymb}
\usepackage{xcolor}
\newtheorem{theorem}{Theorem}
\newtheorem{lemma}{Lemma}
\renewcommand\le\leqslant
\title{{Terminal} Absoluteness of Collapse Forcings}
\author{Cesare Straffelini}
\date{Source: arXiv:2512.16330v2 (CC BY 4.0)}
\begin{document}
\maketitle

\noindent\emph{Source and licence.} {Terminal} Absoluteness of Collapse Forcings. Version arXiv:2512.16330v2 (CC BY 4.0), distributed by arXiv under the Creative Commons Attribution 4.0 International licence, http://creativecommons.org/licenses/by/4.0/, as recorded in the arXiv licence metadata for this exact version and shown on the abstract page https://arxiv.org/abs/2512.16330v2. Full licence notice: \texttt{licenses/arxiv-2512.16330-CC-BY-4.0.txt}.\par\smallskip

\noindent\textit{Editorial scope.} Counted unit: the article's theorem martincons (source0.tex lines 314--317). The article's complete original proof of it is delivered with it (source0.tex lines 320--339, one \texttt{proof} environment of 20 lines). No mathematics is written, repaired or re-derived: the statement and the proof are the article's own lines, in the article's own notation, and no retained line is substituted or re-flowed.  Retained uncounted as the source-local context the proof needs: lines 79--81; lines 85--88; lines 307--310.  Nothing was omitted from any retained span, so the package carries no omission record.  No retained line contains a citation, so the wrapper carries no bibliography and no bibliography entry.  No retained line contains a figure environment, an image-inclusion command or drawing code, so no image is delivered and the package declares no asset dependency.  Editorial formatting only, in addition: an AMS \texttt{amsart} preamble that restates the article's own declarations and macros used by the retained lines, namely the environments \texttt{theorem}, \texttt{lemma} on separate counters, together with the standard packages those lines need.  The retained environments print in this document as Theorem 1, Lemma 1 in order of appearance, whereas the article prints the same items with the numbering of its own full document.  The complete native source of the article as distributed in the arXiv e-print for this exact version is bundled unchanged as \texttt{sources/191234/source0.tex}.
\begin{theorem}[Shoenfield]\label{Shoenfield}
    $\boldsymbol\Sigma^1_2$-$\mathbb P$-Absoluteness holds, for any forcing notion $\mathbb P$.
\end{theorem}

\begin{lemma}\label{upwards}
    Assume $\boldsymbol\Sigma^1_n$-$\mathbb P$-Absoluteness holds, for some $\mathbb P$. For any $\Sigma^1_{n+1}$-formula $\varphi(x_1,\ldots,x_k)$ and any choice of parameters $a_1\subseteq\omega$, $\ldots$, $a_k\subseteq\omega$, we have
    \[\varphi(a_1,\ldots,a_k) \to \mathbb P\Vdash \varphi(\check a_1,\ldots,\check a_k).\]
\end{lemma}

\begin{theorem}[Martin]\label{martincountable}
Assume $a^\sharp$ exists for all $a\subseteq\omega$, and let $\psi$ be a $\Pi^1_2$-formula. Then, there are a formula $\varphi$ and $k<\omega$ such that, for each $x\subseteq\omega$, each $u\subseteq\omega$ such that $x\in \mathsf L[u]$, and each choice of indiscernibles $c_0<\cdots<c_{k-1}$ for $\mathsf L[u]$,
\[\psi(x) \Leftrightarrow \mathsf L[u]\models \varphi(x,c_0,\ldots,c_{k-1}).\]
\end{theorem}

\begin{theorem}\label{martincons}
If $a^\sharp$ exists for all $a\subseteq\lambda$, then $ \boldsymbol\Sigma^1_3\text{-}\mathrm{Coll}(\omega,\lambda)\text{-Absoluteness}$ holds.%, for $\lambda$ a cardinal. If $|\mathbb P|\le \lambda$, then
%\[\mathbb P\Vdash \boldsymbol\Sigma^1_3\text{-}\mathrm{Coll}(\omega,\check\lambda)\text{-Absoluteness}.\]
\end{theorem}

\begin{proof}
Let $G\subseteq\mathrm{Coll}(\omega,\lambda)$ be $\mathsf V$-generic. Thanks to Lemma \ref{upwards} and to Shoenfield's Theorem \ref{Shoenfield}, we only need to prove downward absoluteness for $\boldsymbol\Sigma^1_3$-statements.\medskip

Thus, let $\psi$ be a $\Pi^1_2$-formula and assume that
\[\mathsf V[G]\models \exists x\subseteq\omega\ \psi(x,b)\]
where $b\in\mathsf V$ is a real parameter. Let $w\in\mathsf V[G]$ be a witness for the formula in $\mathsf V[G]$. Our goal is to show that such a witness also exists in $\mathsf V$.\medskip

Let $\tau$ be a $\mathrm{Coll}(\omega,\lambda)$-name for $w$. Since $\tau$ can be coded into a subset of $\lambda$ in $\mathsf V$, in particular $\tau^\sharp$ exists. Since $b$ is a real, $b^\sharp$ also exists, thus $\langle b,\tau\rangle^\sharp$ exists.\medskip

Let $\varphi$ and $k<\omega$ be the ones prescribed by Martin's Theorem \ref{martincountable} for $\psi$. We are going to assume that $k=0$; the case with more indiscernibles is analogous, but more cumbersome in writing. By hypothesis, since $w\in \mathsf L[b,\tau][G]$,
\[\mathsf L[b,\tau][G] \models \varphi(w,b).\]
This, in turn, implies that (since $\mathrm{Coll}(\omega,\lambda)\in\mathsf L[b,\tau]$)
\[\mathsf L[b,\tau] \models \mathrm{Coll}(\omega,\lambda)\Vdash \varphi(\tau,b).\]
Now, let $\mathsf N$ be a countable elementary substructure of $\mathsf L[b,\tau]$ such that $b\in\mathsf N$, $\tau\in \mathsf N$ and $\lambda\in\mathsf N$. Let $\mathsf M$ be the Mostowski collapse of $\mathsf N$, and $\pi:\mathsf N\to\mathsf M$ be the collapse map. Notice that $\pi(b)=b$ since $b\subseteq\omega$ is a real. Then, by elementarity,
\[\mathsf M\models \mathrm{Coll}(\omega,\pi(\lambda))\Vdash \varphi(\pi(\tau),b).\]
Since $\mathsf M$ is countable, in $\mathsf V$ there is a $\mathrm{Coll}(\omega,\pi(\lambda))$-generic filter $H$ over $\mathsf M$. Thus, \[\mathsf M[H]\models \varphi(\imath_G(\pi(\tau)),b).\]
Notice now that $\mathsf L[b,\tau]$ is a model of the sentence ``there is $u$ such that $\mathsf V=\mathsf L[u]$''. This is expressible as a first-order statement, hence $\mathsf M$ is also a model of it. So, we can apply Martin's Theorem \ref{martincountable} again, finding that
\[\mathsf V\models \psi(\imath_G(\pi(\tau)),b),\]
hence we found a witness $w':=\imath_G(\pi(\tau))\in\mathsf V$, as desired.
\end{proof}

\end{document}
